% --- Archon physics formalization source begin ---
% archon:physics
% archon:chemistry
% archon:covers IChO2026Problems/problem_icho_2026_t6_a3.lean
% archon:source-report reports/icho_2026/problem_icho_2026_t6_a3.source.json
% archon:problem-id icho_2026_t6
% archon:part-id T6-A3

\section{Icho Chemistry problem icho\_2026\_t6\_a3}

\paragraph{Problem source.}
\#\# Source
58th International Chemistry Olympiad, Tashkent, Uzbekistan, 2026.
Official-materials index: https://www.icho2026.uz/problems
English problem PDF: ../icho\_2026\_source/raw/theory\_problem.pdf
English solution/rubric PDF: ../icho\_2026\_source/raw/theory\_solution.pdf
Problem PDF page 53; printed local question page 2; solution starts on PDF page 50.
The page PNG is primary visual evidence for formulas, structures, tables, and figures.

\#\# T6. Carbon Nanorings
T6. Carbon Nanorings 7\% In 2019, C18, the first example of a new class of carbon allotropes called cyclo[n]carbons (C𝑛) was detected. Cyclocarbons are monocyclic, all-carbon compounds with high strain energy and low stability. They can be generated and analysed by bond-resolved Atomic Force Microscopy (AFM). Cyclocarbons have interesting aromatic properties based on their π systems. C13 has triplet (3C13, spin S = 1), and singlet (1C13, spin S = 0) forms. 6.1 Based on Hückel’s rule, fill in the table with the number of distinct aromatic (A) and anti-aromatic (AA) 𝜋-systems present in each of C18, C16, 3C13, and 1C13. Fill in all blanks. Later, cyclo[n]carbons C10, C12, C14, C20, and C26 were generated by a combination of dehalogenation and retro-Bergman (1) reactions on a surface. AFM images taken during the synthesis of C14 are shown below: 6.2 Draw structures B-D. Note they may contain one or more unpaired electrons. The structure of A is given as an example.

\#\# Current subquestion 6.3
The voltages applied during the AFM-mediated synthesis of C𝑛may vary. Using the given C–X bond energies, choose the possible halogen(s) in the halogenated reagent for C18 synthesis via a retro-Bergman reaction with a voltage of 2.5 V acting on the electrons. Bond (C–X) Bond Dissociation Energy / kJ mol−1 C–F 467 C–Cl 346 C–Br 290 C–I 228

\#\# Dataset metadata
IChO 2026; T6; raw rubric points=2; kind=classification; formalization\_ready=true.

\paragraph{Shared problem context.}
T6. Carbon Nanorings 7\% In 2019, C18, the first example of a new class of carbon allotropes called cyclo[n]carbons (C𝑛) was detected. Cyclocarbons are monocyclic, all-carbon compounds with high strain energy and low stability. They can be generated and analysed by bond-resolved Atomic Force Microscopy (AFM). Cyclocarbons have interesting aromatic properties based on their π systems. C13 has triplet (3C13, spin S = 1), and singlet (1C13, spin S = 0) forms. 6.1 Based on Hückel’s rule, fill in the table with the number of distinct aromatic (A) and anti-aromatic (AA) 𝜋-systems present in each of C18, C16, 3C13, and 1C13. Fill in all blanks. Later, cyclo[n]carbons C10, C12, C14, C20, and C26 were generated by a combination of dehalogenation and retro-Bergman (1) reactions on a surface. AFM images taken during the synthesis of C14 are shown below: 6.2 Draw structures B-D. Note they may contain one or more unpaired electrons. The structure of A is given as an example.

\paragraph{Current subquestion.}
The voltages applied during the AFM-mediated synthesis of C𝑛may vary. Using the given C–X bond energies, choose the possible halogen(s) in the halogenated reagent for C18 synthesis via a retro-Bergman reaction with a voltage of 2.5 V acting on the electrons. Bond (C–X) Bond Dissociation Energy / kJ mol−1 C–F 467 C–Cl 346 C–Br 290 C–I 228

\paragraph{Recorded answer/context.}
First, we need to convert eV to kJ mol−1: 𝐸𝑔𝑖𝑣𝑒𝑛[kJ mol−1] = 2.5 × (1.602 × 10−19 × 6.022 × 1023)/1000 kJ mol−1 (1) 𝐸𝑔𝑖𝑣𝑒𝑛= 2.5 ⋅96.485 kJ mol−1 = 241.2 kJ mol−1 (2) The only candidate with 𝐵𝐷𝐸≤𝐸𝑔𝑖𝑣𝑒𝑛is the C–I bond. X = I 1 point for unit transformation, 1 point for correct halogen choice. (2 pts awarded if I is chosen with no working) 0 pts for other halogens or more than one halogen chosen. 2.0 pt

\paragraph{Figure/image paths.}
\begin{itemize}
\item ../icho\_2026\_source/image/T6\_page-2.png
\item ../icho\_2026\_source/image/T6\_page-1.png
\end{itemize}

\paragraph{Formalization target.}
create a compiling Lean file with sorry bodies at `IChO2026Problems/problem\_icho\_2026\_t6\_a3.lean`.
The Lean declarations must preserve every mathematical object, quantifier, hypothesis, side condition, approximation/error guarantee, and requested conclusion from the source statement.
Use Mathlib/Physlib/Chemistry names found through LeanExplore where available. Prefer the configured benchmark Base library; introduce a local definition only when the source genuinely requires an object that the environment does not expose.

\begin{theorem}\leanok
[Icho Chemistry formalization target]
\label{thm:physics:icho_2026_t6_a3:target}
\lean{IChO2026Problems.T6A3.possible_halogen_for_C18_retroBergman_at_2_5_volts}
The available electrochemical energy at 2.5 V selects the possible halogen for
the C18 retro-Bergman reaction.
\end{theorem}
\begin{proof}

\leanok
The supplied conversion gives $2.5\times96.485=241.2125\,
\mathrm{kJ\,mol^{-1}}$.  Compare this with the four tabulated dissociation
energies: $467$, $346$, and $290$ for C--F, C--Cl, and C--Br exceed the
available energy, while the C--I value $228$ does not.  Under the stated C18
retro-Bergman context, energetic cleavability is therefore equivalent to the
halogen being iodine.
\end{proof}
% --- Archon physics formalization source end ---
